\section{Planteamiento del Problema}

El mundo actual se distingue por la aceleración con la que se vive, lo que no
solo provoca un avance tecnológico cada vez mayor, sino que además trae consigo
la aparición de nuevos productos cuyo objetivo es resolver ``necesidades'' de la
sociedad; sin embargo, muchas de estas necesidades no son reales, sino que —en
palabras de \textcite[p.~347]{martinez2016}— son generadas por la maquinaria
publicitaria para inducir al consumidor a adquirir nuevos modelos antes de que
sea realmente necesario, manteniéndolo en un ciclo de consumo desmedido que
incrementa la demanda y, con ella, el uso de recursos.

Tomando en cuenta lo anterior, el consumidor promedio se ve envuelto en una
sobrecarga de información en donde lo único que conoce son las necesidades
impuestas por el capitalismo moderno; deja de estar interesado en lo que
conlleva la creación de un producto, buscando lo que aparentemente cumple con lo
que el consumidor busca, pues las campañas de marketing muestran únicamente lo
que las empresas quieren que el consumidor vea, lo que no siempre concuerda con
la realidad, generando confusión sobre la calidad técnica de un producto.
Además, la cantidad de información que un producto suele tener en sus fichas
técnicas suele provocar que, al vivir en una época de inmediatez, el consumidor
no se permita el tiempo de leer —o informarse— a profundidad. 

Además de los problemas que la inmediatez del mundo actual general, el
confinamiento debido a la pandemia COVID-19 tuvo sus propias consecuencias
dentro de la problemática del consumismo, donde muchas personas encontraron
refugio para sobrellevar el aislamiento a través de las compras en línea. Este
comportamiento se vio reflejado directamente en la cantidad de usuarios nuevos
en aplicaciones e-commerce; solo en 2019, el uso de los medios digitales
incremento un 24\% en relación con el 2017 en México, con un valor de 491
millones de pesos \parencite[p.~123]{bejar2023}. El hecho de comprar en línea
resultó una actividad sencilla para el consumidor, sin requerir ningún tipo de
esfuerzo, por lo que millones de personas comenzaron a unirse a este tipo de
actividad sin realmente tener un objetivo en sus compras. 

La problemática del consumismo tiene múltiples causas; sin embargo, la actual
investigación estará enfocada principalmente en dos estrategias empleadas por
las empresas para mantenerlo. El uso del llamado greenwashing —la práctica de
presentar productos y campañas de marketing como ecológicos o sostenibles sin
serlo realmente— genera una percepción de responsabilidad ambiental que es, en
esencia, engañosa. Así como la sobreabundancia de información que confunde al
consumidor, desalentando cualquier interés por conocer a fondo la composición o
el impacto real de lo que adquiere. Ambas estrategias, combinadas con la
obsolescencia programada —el diseño deliberado de productos con una vida útil
limitada—, garantizan que el ciclo del consumismo se mantenga activo de forma
ininterrumpida, sin que el consumidor encuentre un momento de pausa para
cuestionar sus hábitos de compra. 

Con base en lo anterior, la presente investigación busca responder a la
siguiente pregunta: \noindent\textit{¿Cómo puede una aplicación móvil de escaneo
de código de barras contribuir a reducir la asimetría de información en la
compra de productos electrónicos (smartphones y laptops) y fomentar un consumo
más informado y sostenible?}